\documentclass[a4paper,10pt]{report}\input{../0.Préambule/Préambule_cours_prof}\begin{document}

\newpage
\pagestyle{empty}
\section*{Statistiques - Utilisation de dés}

\vspace{-3mm}
\setcounter{cexer}{0}
\begin{exer1}

\begin{enumerate}[font=\color{exer1}]
\item Lancer un dé à $6$ faces $20$ fois chacun, note le résultat obtenu.

\renewcommand{\arraystretch}{1.3}
\begin{tabularx}{\linewidth}{|>{\columncolor{exer1!20}}p{1.5cm}|*{20}{>{\centering \arraybackslash}X|}}
\hline
 & & & & & & & & & & & & & & & & & & & \\\hline
 & & & & & & & & & & & & & & & & & & & \\\hline
\end{tabularx}

\item Regrouper les résultats obtenus dans le tableau suivant :

\renewcommand{\arraystretch}{1.3}
\begin{tabularx}{\linewidth}{|>{\columncolor{exer1!20}}p{3cm}|*{6}{>{\centering \arraybackslash}X|}}
\hline
Résultat du dé  & $1$ & $2$ & $3$ & $4$ & $5$ & $6$ \\\hline
Effectif & & & & & & \\\hline
Fréquence & & & & & & \\\hline
Fréquence en $\%$ & & & & & & \\\hline
\end{tabularx}
\item Complète le tableau en calculant la fréquence et la fréquence en $\%$.
\end{enumerate}
\end{exer1}

\begin{exer2}

\begin{enumerate}[font=\color{exer2}]
\item Lancer 20 fois chacun deux dés à 6 faces, additionne la valeur des deux dés, note le résultat.

\renewcommand{\arraystretch}{1.3}
\begin{tabularx}{\linewidth}{|>{\columncolor{exer2!20}}p{1.5cm}|*{20}{>{\centering \arraybackslash}X|}}
\hline
 & & & & & & & & & & & & & & & & & & & \\\hline
 & & & & & & & & & & & & & & & & & & & \\\hline
\end{tabularx}
\item Présente les résultats dans un tableau comme pour l'exercice précédent.

\renewcommand{\arraystretch}{1.3}
\begin{tabularx}{\linewidth}{|>{\columncolor{exer2!20}}p{3cm}|*{1}{>{\centering \arraybackslash}X|}}
\hline
Résultat du dé  & \\\hline
Effectif & \\\hline
Fréquence & \\\hline
Fréquence en $\%$ &\\\hline
\end{tabularx}
\end{enumerate}
\end{exer2}

\vspace{5mm}
\section*{Statistiques - Utilisation de dés}

\vspace{-3mm}
\setcounter{cexer}{0}
\begin{exer1}

\begin{enumerate}[font=\color{exer1}]
\item Lancer un dé à $6$ faces $20$ fois chacun, note le résultat obtenu.

\renewcommand{\arraystretch}{1.3}
\begin{tabularx}{\linewidth}{|>{\columncolor{exer1!20}}p{1.5cm}|*{20}{>{\centering \arraybackslash}X|}}
\hline
 & & & & & & & & & & & & & & & & & & & \\\hline
 & & & & & & & & & & & & & & & & & & & \\\hline
\end{tabularx}

\item Regrouper les résultats obtenus dans le tableau suivant :

\renewcommand{\arraystretch}{1.3}
\begin{tabularx}{\linewidth}{|>{\columncolor{exer1!20}}p{3cm}|*{6}{>{\centering \arraybackslash}X|}}
\hline
Résultat du dé  & $1$ & $2$ & $3$ & $4$ & $5$ & $6$ \\\hline
Effectif & & & & & & \\\hline
Fréquence & & & & & & \\\hline
Fréquence en $\%$ & & & & & & \\\hline
\end{tabularx}
\item Complète le tableau en calculant la fréquence et la fréquence en $\%$.
\end{enumerate}
\end{exer1}

\begin{exer2}

\begin{enumerate}[font=\color{exer2}]
\item Lancer 20 fois chacun deux dés à 6 faces, additionne la valeur des deux dés, note le résultat.

\renewcommand{\arraystretch}{1.3}
\begin{tabularx}{\linewidth}{|>{\columncolor{exer2!20}}p{1.5cm}|*{20}{>{\centering \arraybackslash}X|}}
\hline
 & & & & & & & & & & & & & & & & & & & \\\hline
 & & & & & & & & & & & & & & & & & & & \\\hline
\end{tabularx}
\item Présente les résultats dans un tableau comme pour l'exercice précédent.

\renewcommand{\arraystretch}{1.3}
\begin{tabularx}{\linewidth}{|>{\columncolor{exer2!20}}p{3cm}|*{1}{>{\centering \arraybackslash}X|}}
\hline
Résultat du dé  & \\\hline
Effectif & \\\hline
Fréquence & \\\hline
Fréquence en $\%$ &\\\hline
\end{tabularx}
\end{enumerate}
\end{exer2}

\end{document}